\section{Marco Teórico}
\subsection{Localización en interiores}
{\justify La localización en interiores, también conocida como Indoor Positioning, es un área de estudio dentro de la computación ubicua que busca determinar la posición de personas u objetos dentro de espacios cerrados, donde las señales satelitales como el GPS pierden cobertura o presentan alta inexactitud. A diferencia de los sistemas de localización en exteriores, la implementación de soluciones en interiores implica múltiples retos técnicos, como la interferencia de estructuras físicas (paredes, techos, mobiliario), el comportamiento impredecible de las señales de radiofrecuencia por efectos como la reflexión, refracción y difracción, así como la variabilidad en el entorno.

Para suplir estas limitaciones, se han explorado diferentes tecnologías de localización alternativas, entre ellas Wi-Fi, Bluetooth Low Energy (BLE), RFID, visión por computadora y Ultra-Wideband (UWB). Cada tecnología ofrece un compromiso distinto entre precisión, consumo energético, costo y complejidad de implementación. En entornos que requieren alta precisión —como hospitales, bodegas o, en este caso, instalaciones universitarias complejas—, el UWB se destaca por ofrecer ventajas clave frente a sus competidores, sobre todo en términos de exactitud, latencia y robustez frente al ruido ambiental.}


\subsection{Tecnología Ultra-Wideband (UWB)}
{\justify La tecnología Ultra-Wideband (UWB) se basa en el uso de pulsos de muy corta duración que se propagan en un espectro extremadamente amplio (mayor a 500 MHz). Esta característica permite transmitir datos con alta precisión temporal, lo que a su vez facilita la medición del tiempo exacto que tarda una señal en desplazarse entre dos dispositivos. A través de estos tiempos de vuelo (Time of Flight, ToF), un receptor puede estimar con mucha precisión la distancia a un emisor, sin depender de la intensidad de la señal (RSSI), como ocurre en tecnologías menos precisas como BLE o Wi-Fi.

En aplicaciones de localización, los dispositivos UWB se colocan en el entorno como anclas o beacons, que transmiten pulsos periódicamente. Un dispositivo móvil equipado con un receptor UWB puede calcular su distancia a múltiples beacons y usar dicha información para estimar su posición. Esta técnica, que combina alta precisión (en el orden de decímetros o incluso centímetros) con baja latencia, es ideal para aplicaciones donde se requiere posicionamiento en tiempo real con gran fidelidad espacial, como la navegación asistida con realidad aumentada. En este proyecto se utilizan sensores Estimote UWB Beacons, diseñados específicamente para tareas de localización en interiores, con soporte en plataformas móviles como iOS.}


\subsection{Trilateración}
{\justify La trilateración es un método geométrico fundamental en sistemas de localización que permite calcular la posición de un punto desconocido en el espacio utilizando las distancias conocidas a tres o más puntos de referencia cuyas ubicaciones son fijas. En el caso de localización en dos dimensiones, si se conocen las coordenadas de tres sensores (beacons) y las distancias del dispositivo móvil a cada uno de ellos, es posible determinar su posición mediante la intersección de los tres círculos generados por esas distancias. En la práctica, debido a errores de medición, estas circunferencias rara vez se intersectan en un solo punto, por lo que se recurre a métodos de optimización y regresión para encontrar el punto más probable.

Cuando se disponen más de tres sensores, se puede aplicar trilateración redundante, técnica que permite mejorar la estabilidad de la estimación al reducir el impacto de errores individuales en las mediciones. La trilateración es la base matemática de muchos sistemas de localización, incluidos GPS, y en este proyecto, se adapta al contexto del campus universitario, utilizando coordenadas cartesianas en un plano 2D para determinar la ubicación del usuario en tiempo real.}

\subsection{ Filtro de Kalman}
{\justify El filtro de Kalman es un algoritmo recursivo de estimación desarrollado por Rudolf E. Kálmán en la década de 1960, utilizado para predecir el estado de un sistema dinámico y corregir esa predicción a partir de mediciones ruidosas. Su aplicación ha sido clave en campos como la navegación inercial, la robótica, el seguimiento de objetos y los sistemas de control, debido a su capacidad para integrar múltiples fuentes de información en presencia de incertidumbre.

En el contexto de este proyecto, el filtro de Kalman se emplea para suavizar la trayectoria estimada del usuario dentro del campus, combinando las posiciones calculadas mediante trilateración con los datos del acelerómetro del dispositivo móvil. Esta fusión sensorial permite reducir el jitter, es decir, las pequeñas oscilaciones aleatorias en la posición del usuario, y mejorar la estabilidad de la trayectoria percibida. Aunque el sistema no busca una precisión milimétrica absoluta, el objetivo es proporcionar una experiencia de navegación fluida, donde la posición calculada varíe de forma coherente con el movimiento del usuario.}

\subsection{Acelerómetro}
{\justify El acelerómetro es un sensor electrónico capaz de medir la aceleración de un cuerpo en los tres ejes del espacio (x, y, z), tanto por efecto del movimiento como por la fuerza gravitacional. En los dispositivos móviles modernos, los acelerómetros permiten funciones como la detección de orientación, el conteo de pasos, y la navegación sin GPS (dead reckoning).

En este proyecto, el acelerómetro se utiliza como complemento a la trilateración UWB para mejorar la estimación de la posición del usuario. Aunque los acelerómetros presentan el inconveniente de la deriva acumulativa (pequeños errores que se amplifican con el tiempo cuando se integra la aceleración), su uso combinado con los datos de distancia de los sensores UWB, mediante el filtro de Kalman, permite detectar transiciones suaves en el movimiento del usuario, anticipar su dirección y estabilizar la trayectoria. De esta manera, el sistema puede ofrecer una experiencia más coherente y natural en la navegación con realidad aumentada.}

\subsection{Unity y Realidad Aumentada}
{\justify Unity es un motor de desarrollo multiplataforma ampliamente utilizado en la industria de los videojuegos, simulaciones y aplicaciones interactivas. Su compatibilidad con dispositivos móviles y su integración con librerías de realidad aumentada como ARKit (iOS) y ARCore (Android), lo convierten en una herramienta poderosa para el desarrollo de experiencias inmersivas en tiempo real.

En este proyecto, Unity se emplea como el entorno de ejecución principal para la aplicación de recorridos virtuales con realidad aumentada. A través de la cámara del dispositivo móvil, el usuario puede observar su entorno mientras el sistema le sobreimpone elementos gráficos —como flechas tridimensionales— que le indican hacia dónde debe avanzar para completar un recorrido. La interacción entre estos elementos visuales y el sistema de localización precisa permite convertir el espacio físico del campus en una experiencia guiada interactiva, intuitiva y accesible.}

\subsection{Plugins Nativos para Unity (iOS)}
{\justify Debido a que Unity no ofrece acceso directo a todas las funcionalidades del sistema operativo, como sensores especializados o bibliotecas nativas, es común extender su funcionalidad mediante plugins nativos. Un plugin nativo permite conectar el código de Unity con código escrito en lenguajes como Swift (para iOS) o Kotlin (para Android), habilitando el uso de SDKs que solo están disponibles de forma nativa.

En el caso de este proyecto, se desarrolla un plugin en Swift que se encarga de establecer la comunicación con los sensores Estimote UWB, procesar las mediciones de distancia, acceder a los datos del acelerómetro mediante el framework CoreMotion, realizar los cálculos de trilateración y aplicar el filtro de Kalman. Este plugin se comunica con Unity a través de una interfaz de funciones expuestas, enviando las coordenadas calculadas para ser usadas por el motor gráfico en la navegación. Esta arquitectura modular garantiza que el sistema pueda mantenerse y extenderse fácilmente, incluyendo la posibilidad futura de implementar un plugin análogo para Android.}
\newpage